\section{Further research questions and a programme of exact identities}
\label{research:section}

These questions begin at the boundaries of the proved results. They
ask for exact formulas, functional relations, or finite certificates;
none is a conjectural assertion presented as a theorem. They are ordered
roughly by proximity to computations already available in this package.

\subsection{Ordered jets, derivatives, and antiderivatives}

\paragraph{Q1. Differential closure of the depth-four coordinates.}
Equation~\eqref{oq:derivative} already removes $T$ in favor of
$\eta'(a)$, ordinary Hurwitz jets, and the convergent value
$Z_2(2,1;a)$. Derive analogous argument-derivative formulas for
$\rho$ and $\kappa$, then use stuffle and partial fraction identities
to reduce the resulting depth-three and mixed-weight sums. A useful
first result would be a finite differential system with all additive
normalizations fixed by the shift laws. Determine the smallest
\emph{proved spanning} system before asking for arithmetic minimality.

\paragraph{Q2. Two or more blocks of equal slopes at arbitrary depth.}
The all-equal and one-distinguished-slope families have Gamma
generators. At depth four, the fully independent formula identifies
the first coefficients that such generators miss. Seek an exact
two-block generator for slopes
$(p,\ldots,p,q,\ldots,q)$, with the block lengths arbitrary. The target
is a finite collection of lower-depth generating germs independent of
the block lengths, or a proof that a proposed finite collection fails.
The depth-four normal form supplies a mandatory first consistency
specialization.

\paragraph{Q3. Depth five with an explicit coordinate ledger.}
At depth five the finite part needs the cubic jet of $R_2$, the
quadratic jet of $R_3$, and the linear jet of $R_4$. Compute the exact
regularized stuffle constraints on these jets and construct an analogue
of the three-direction reconstruction. The deliverable should include
a convergent representative for every retained coordinate, its shift
recurrence, and a symbolic certificate for each claimed cancellation.
This is a finite linear problem after the analytic regular germs have
been fixed; it is not automatically a period-independence problem.

\paragraph{Q4. Rational-shift character combinations.}
For conductors $2,3,4,6$, apply the exact root-of-unity filter to the
convergent representatives of $\eta,\rho,\kappa,T$. Determine which
character-weighted combinations eliminate the residual ordered
coordinates and reduce to ordinary Hurwitz or Dirichlet $L$ jets.
The shift identities alone give translations, not distribution.
A successful theorem must show the cancellation algebraically and
retain every logarithmic scale term.

\paragraph{Q5. Values on a polar hyperplane.}
The transverse-ray formula assumes all prefix slopes are nonzero.
For a family lying in one prefix hyperplane, compare three explicit
prescriptions: take a finite part in a normal variable first, remove
the ordered poles through $R_d$ first, or approach along a specified
two-parameter surface. Compute their exact differences at depth four
and determine which prescriptions preserve a stated stuffle identity.
No value is defined by substituting a zero denominator into the ray
formula.

\subsection{Coincident products and Gamma primitives}

\paragraph{Q6. The first two-unreflected Stieltjes closure problem.}
Reflection evaluates $\FP\int\psi^2K^{2h+1}$ in this article, while
one arbitrary Stieltjes derivative times a reflected polynomial always
closes. Classify the combinations of
\[
 \FP\int_0^1\gamma_m^{(p)}(x)\gamma_n^{(q)}(x)K(x)^r\,dx
\]
whose completed higher-depth coordinate cancels. Begin with
$m+n\le2$ and total argument derivative order at most two. The output
should be an exact kernel of a finite coordinate map, with a proof of
each cancellation and fixed endpoint coordinates.

\paragraph{Q7. Evaluate the generic cubic Tornheim coordinate.}
The cubic difference \eqref{ct:cubic-difference} reduces a mixed
moment without determining the common unreflected cubic value. Seek a
functional relation for the diagonal Tornheim third derivative in the
incoming cubic formula, or a normalized Barnes-type representation
whose derivative can be proved to equal it. This must be a new analytic
relation, not another rearrangement of an equivalent cubic identity.
The existing Mellin representation is a viable starting point.

\paragraph{Q8. A reflected generating function beyond polynomials.}
Polynomial closure is finite because the cotangent derivative
polynomials form a triangular basis. Determine a nonpolynomial family
$P_\lambda(K)$ for which the completed integral has a holomorphic
parameter generator and controlled coefficientwise endpoint completion.
Trigonometric beta integrals suggest rational or complex-power
families, but their changing endpoint exponents introduce new
resonances. A useful theorem would identify the full local subtraction
before any expansion in $\lambda$.

\paragraph{Q9. Multiple Gamma primitives with fixed constants.}
The log-Gamma moments here are first argument primitives of digamma
moments. Construct analogous Barnes-$G$ or higher-Gamma identities
for repeated primitives, with the integration constant fixed at an
explicit positive point and compatibility checked against the
functional equation. The endpoint constant terms must be included
at every integration by parts. A first target is a closed generator
for the even cotangent moments at the next primitive level.

\subsection{Harmonic powers and spectral regular coefficients}

\paragraph{Q10. Compact all-degree formulas for centered values.}
The finite word recursion proves that every
$F_p(-2r,1/2)$ belongs to the ordinary multiple-zeta algebra.
Seek a bivariate generator in $p$ and $r$ that improves on enumerating
all compositions of $p$. Even a closed formula for the entire sequence
$F_p(-2,1/2)$ would extend the explicit table. State the exact MZV
coordinate space used at each weight rather than assume every value
reduces to products of ordinary zeta values.

\paragraph{Q11. Generalized harmonic products and their regular shifts.}
For a fixed polynomial in $H_n^{(r)}$, $r\ge1$, the same Mellin
calculus produces polynomial finite parts. Which numerator combinations
make every nonpositive even principal part vanish at $a=1/2$?
The raw-power parity proof does not apply automatically when
$H_n^{(r)}$ with $r>1$ is present. Compute the finite Bernoulli
conditions at successive orders and seek a necessary and sufficient
generating criterion. This gives exact cancellation identities rather
than inequalities.

\paragraph{Q12. Cancel the first global Mellin remainder.}
The first regular spectral coefficient at $s=-m$ contains global
information absent from its finite part. Find finite combinations of
harmonic-power series whose remainder integrals cancel by an exact
kernel identity. Proposition~\ref{hp:prop:tin-all-log} is a model:
one differential combination has a normally convergent series and an
explicit relation among higher harmonic Stieltjes constants. Generalize
this cancellation to nonpositive centers and generalized harmonic
numerators, with the normalization checked independently.

\subsection{The Gaussian period candidates}

\paragraph{Q13. Identify additional information detected by the separator.}
The inherited functional survives completion of the entire Cayley
ideal. A useful next computation is to generate selected additional
double-shuffle or distribution relations outside the frozen family,
project them by $\Pi_\uparrow$, and identify a row on which the
functional is nonzero. Such a row is a concrete source of information
missing from the present search; it is not by itself a proof of $S_6$.
Continue until either an exact rational proof certificate for the
target is obtained or a new fully specified separating certificate is
proved. Higher-depth intermediate rows are permitted, but need not be
assumed necessary a priori.

\paragraph{Q14. A functional bridge to $S_6$ and the revised $S_8$.}
Seek a parameter-dependent polylogarithm identity, differential system
with enough exact boundary data, or regulator relation whose
specialization is the short candidate. Independently, repeat the
precisely defined formal quotient analysis for the revised weight-nine
target. The weight-seven certificate cannot be extrapolated to it.
Any new numerical relation remains a candidate until such a functional
or exact symbolic proof is supplied.

\paragraph{Q15. Formal verification of the finite-part and ideal mechanisms.}
The local completion lemma and the depth-preserving ideal-intersection
lemma are reusable proof units. Formalize those units in a proof
assistant, then import the finite rational certificate with an exact
integer parser. For the analytic identities, formalize normal
convergence after a finite local subtraction before importing
coefficient formulas. This would turn the current reproducible
computations into a smaller formally checked core without confusing
floating-point diagnostics with theorem proofs.
